% Ad Atticum 3.15 (ad-atticum-03-15) --- English translation
% Translated by Claude (Anthropic) from the Latin (Perseus).
% Style guide: STYLE.md.

\ciceroLetterOpener{CICERO ATTICO salutem}

\ciceroSection{1} On the Ides of August I received four letters sent from you: one in which you take me to task and ask me to be more steady; another in which you say that a freedman of Crassus's told you of my anxiety and the wasting of my body; a third in which you set out the proceedings in the Senate; a fourth on the matter, as you write, in which Varro has confirmed to you what Pompey is minded to do.

\ciceroSection{2} To the first I write this in answer: that I grieve in such a way that I am not in the least deserted by my mind, but I grieve at this very thing --- that, having so firm a mind, I have no men to use it with, no men with whom. For if you, without your grief, do not lack me alone, what do you think of me, who lack both you and everything? And if you, safe yourself, miss me, how do you suppose that I myself miss safety itself? I do not want to recall what I have been stripped of: not only because the things are not unknown to you, but also that I may not myself reopen my grief. This I affirm: that no one has been deprived of such goods, nor fallen into such miseries. Time, moreover, not only does not lighten this grief but increases it. Other griefs are softened by length of years: this one cannot but be daily increased, by the feel of the present misery and the recollection of the life that was. For it is not only my own things, my own people that I miss, but myself. What, after all, am I? But I will not bring myself either to wear down your spirit with complaining, or to bring my hands again and again to my own wounds. As to your clearing of those whom I had written had envied me, Cato among them --- I think that man was so far from that villainy that what I most grieve at is this: that the pretence of others had more weight with me than his good faith. As to your clearing the rest, if they are acceptable to you, they ought to be acceptable to me. But these are things we discuss too late.

\ciceroSection{3} Of Crassus's freedman I think nothing was honestly said. You write that the matter was well handled in the Senate. But what of Curio? Did he not read that speech? Where it has come out from, I do not know. Axius, writing to me of the proceedings of the same day, does not so praise Curio. He can let some things slip; you, certainly, will not have written what was not the case. Varro's conversation makes one expect Caesar. Would that Varro himself would lean his weight into the cause! And surely he will, both of his own accord and with you pressing.

\ciceroSection{4} If fortune ever puts me again in possession of you and of my country, I shall surely bring it about that you, alone of all my friends, shall rejoice the most. The duties and zealous services I owed you, which shone too little before --- one must confess it --- I shall so carry through that you will think me restored to you on the same footing as to my brother and to our children. If I have offended you in anything --- or rather, since I have offended --- pardon me; for I have offended myself far more harshly. I do not write this because I do not know that you have been struck by my disaster with the greatest grief; but surely, if you should have loved and ought to have loved me as much as you do love me and did love me, you would never have let me go without the counsel in which you abounded, nor have let me be persuaded that the law about the colleges, of all things, was useful to us. As it was, you only gave my grief tears, which was an act of love --- just as I myself gave it. What my deserts could have brought about --- that day and night you should be considering what I had to do --- that was passed over, and by my crime, not yours. For if not you alone, but anyone, had been there to call me back, when I was terrified by Pompey's stinted answer, from a most disgraceful counsel --- a thing which only you, or chiefly you, could have done --- either I should have fallen with honour, or we should be living today as victors. Forgive me here. For I am accusing myself far more, and you only as a second self and at the same time a partner in my own fault. And if I am restored, even our offence will look the smaller; and you and I, certainly, shall love each other from now on by your service alone, since by no service of mine.

\ciceroSection{5} As for what you write of having spoken with Culleo about the \textit{privilegium}: there is something in that, but it is far better to have it abrogated. If no one obstructs, that way is firmer; and if any man shall not allow it to be carried, the same will obstruct a senatorial decree as well. And nothing more is needed than that it be abrogated --- for the earlier law did me no harm. Had I been willing, when it was promulgated, either to praise it, or to neglect it as it ought to have been neglected, it could have done me no hurt at all. Here, for the first time, my own judgement failed me --- and not only failed but did me harm. Blind, I say blind, we were in changing our dress, in entreating the people. Once it had not begun to be moved against me by name, to do these things was destruction. But I keep going back to the past --- and for this one reason: that, if anything is to be moved, you may not touch the law in which there are many popular provisions.

\ciceroSection{6} But it is foolish of me to be telling you what you should do, or how. Would only that something were being done! In which very thing your letters hide much, I think, lest I be too violently shaken by despair. For what do you see can be done, or how? Through the Senate? But you yourself wrote that Clodius had fixed a clause of the law on the doorpost of the Curia, that nothing should be referred or said about it. How then did Domitius say he would refer it? How was Clodius silent when those men whom you write of were both speaking on the matter and demanding that it be referred? And if it goes through the people, will it be possible without the consent of all the tribunes? What of my goods? what of my house? Can either be restored? Or if not, how shall I myself be? Unless you see all this set in train, into what hope are you calling me? But if there is no hope, what life is there for me? So I am waiting at Thessalonica for the proceedings of the Kalends of August, on the strength of which I shall decide either to take refuge in your fields --- to be away from men I do not wish to see, and to see you, as you write, and to be nearer if anything is being done (and I understand that this is the wish both of you and of my brother Quintus) --- or to go off to Cyzicus.

\ciceroSection{7} Now, Pomponius, since you imparted nothing of your good judgement to my safety --- because either you had decided that there was counsel enough in me alone, or that you owed me nothing further than to be at hand --- and since I, betrayed, drawn on, flung into the snare, neglected all my defences, abandoned and left all Italy when she was already on her feet to defend me, gave myself and mine over to my enemies, with you looking on and silent --- you, who, even if you did not have more talent than I had, were certainly less frightened: if you can, raise up the men brought low, and help us in that; but if everything is blocked, see to that very thing, that we know it, and stop at last either rebuking me or jointly consoling me. If I were accusing your good faith, I would not be the man to commit myself to your roof of all places. I accuse my own madness, in supposing that I was loved by you as much as I would have wished. Had this been so, you would have brought to bear the same good faith, but greater care --- you would surely have held me back as I rushed to my own destruction --- you would not now be undertaking these labours that you take up in our shipwreck.

\ciceroSection{8} Therefore see to it that you write me everything closely examined and explored, and that you want me, as you do, to be somebody --- since what I was, and what I could have been, I cannot now be. Take it that by this letter you have been accused not by me, but I myself by myself. If there are men to whom you think a letter in my name should be sent, I should wish you to write them out and see to their being sent. Sent the fourteenth day before the Kalends of September.
